%% ch06_microgrid_protection.tex
%% Chapter 6: Generator Protection Suite for IBR-Affected Networks
%% Source TRs: TR-56 (87G DFIG), TR-58 (Relay 78 EAC), TR-60 (Relay 40 LOE),
%%             TR-61 (Relay 64G), TR-62 (87L-PV), TR-63 (Relay 81 IFDI)
%% All six TRs complete April 2026.
%% Target journals: IEEE TEC / TPWRS / TPWRD / TIE / TSG
%% ============================================================
\chapter{Generator Protection Suite for IBR-Affected Networks}
\label{ch:generator_suite}
\label{ch:microgrid}% compatibility alias for prior cross-references

% ── Chapter contributions ─────────────────────────────────────────────────────
\begin{itemize}
  \item[\textbf{C9.}] Physics-reduction proposition (Proposition~C9): replacing
    one collinear or derivable parameter reduces the estimator condition number
    $\kappa_n$ by $\geq 10\times$, proved via the SVD interlacing theorem.
    Applied across all six generator-suite zone models.

  \item[\textbf{C10.}] DFIG 10-parameter piecewise-ODE model with CUSUM crowbar
    detection (TR-56, Section~\ref{sec:ch6_87g}): shared stator transient time
    constant $\tau_s$ for DC and natural-response components reduces parameter
    count from 11 to 10 and gives $\kappa_n \in [13, 67]$ vs.\ $\kappa_n \in
    [3, 36]$ for the SG model.
    CCR\,=\,100\% across all slip groups; ANDES WTDTA1 6/6 validated.

  \item[\textbf{C11.}] Two-parameter solar PV zone model with $k_{\rm ibr}$-driven
    DDR model selector (TR-62, Section~\ref{sec:ch6_87l_pv}):
    analytically proven $\kappa_n = 1$, reducing 87L trip time for PV-fed lines
    from $\geq 40\,\text{ms}$ to $20\,\text{ms}$;
    $P_D = 0.9984$, $P_{FA} = 0.0000$ (5\,000-trial MC).

  \item[\textbf{C12.}] IFDI-based Relay~81 for zero-inertia GFM networks
    (TR-63, Section~\ref{sec:ch6_relay81}): 6-parameter model, two-pass LM,
    IFDI trip threshold 1.15, achieving $P_D = 0.9947$, $P_{FA} = 0.0032$.

  \item[\textbf{C13.}] Conservative CUEP lower bound for Relay~78 (TR-58,
    Proposition~1, Section~\ref{sec:ch6_relay78}):
    $\delta_{\rm LB} = \pi - \arcsin(P_m / P_{\rm maxeff})$,
    proved $\leq \delta_u$ always (0 violations / 2\,000 MC trials),
    providing a formal IBR-corrected out-of-step blinder setting.

  \item[\textbf{C14.}] Relay~40 GFM admittance correction (TR-60,
    Section~\ref{sec:ch6_relay40}): online correction $Y_{\rm corr} =
    Y_{\rm total} - \widehat{Y}_{\rm GFM}$ eliminates both nuisance trip and
    masked LOE failure modes; $P_D: 0.568 \to 1.000$, $P_{FA}: 0.038 \to 0.000$.

  \item[\textbf{C15.}] Relay~64G optimal sub-harmonic injection (TR-61,
    Section~\ref{sec:ch6_relay64g}): adaptive $f^* \in [25, 36]\,\text{Hz}$
    maximises balanced security-dependability margin; $P_D: 0.813 \to 1.000$,
    $P_{FA}: 0.819 \to 0.000$ vs.\ conventional 20\,Hz injection.
\end{itemize}

% ──────────────────────────────────────────────────────────────────────────────
\section{Introduction}
\label{sec:ch6_intro}

Five relay functions form the last line of defence for synchronous machines and
their IBR equivalents: Relay~40 (loss of excitation), Relay~64G (stator earth
fault), Relay~78 (out-of-step), Relay~81 (frequency deviation), and the
generator differential (87G).
All five were designed for synchronous machines in SG-dominated grids, and all
five exhibit systematic failures when the generation mix includes grid-following
(GFL) or grid-forming (GFM) inverter-based resources.
Section~\ref{sec:ch6_87l_pv} addresses a sixth function: the 87L line
differential relay on lines served primarily by solar PV inverters, which
requires a dedicated two-parameter zone model.

\begin{description}
  \item[87G / 87L (DFIG):] DFIG fault current contains a slip-frequency natural
    response at $f_s = s\,f_0$ ($s \in [-0.30, +0.30]$) after crowbar firing.
    The standard 4-parameter SG model cannot fit this component, causing the
    relay to misclassify fault current as load flow or to trip spuriously on the
    natural response transient.
    (Section~\ref{sec:ch6_87g}, TR-56.)

  \item[87L (PV):] Grid-following PV inverters produce fault current with no DC
    exponential component.  The 4-parameter SG model wastes two degrees of
    freedom fitting noise, delaying trip confirmation by 20--40\,ms and elevating
    the effective detection threshold.
    (Section~\ref{sec:ch6_87l_pv}, TR-62.)

  \item[Relay~81 (frequency):] In zero-inertia GFM islands, frequency dynamics
    are governed by virtual droop with time constant $\tau_f \in [0.05, 1.0]\,\text{s}$
    — a 20$\times$ range that makes fixed-threshold ROCOF relays unreliable.
    (Section~\ref{sec:ch6_relay81}, TR-63.)

  \item[Relay~40 (LOE):] GFM IBR with virtual-synchronous-machine control alters
    the apparent admittance trajectory in ways that both cause nuisance trips
    (reactive-absorption mode) and mask genuine LOE events (reactive-supply mode).
    (Section~\ref{sec:ch6_relay40}, TR-60.)

  \item[Relay~78 (OOS):] The equal-area criterion blinder assumes all
    synchronising power is supplied by the SG.  For $\geq 30\%$ GFL penetration,
    the standard blinder is non-conservative: it lies outside the true critical
    unstable equilibrium point (CUEP).
    (Section~\ref{sec:ch6_relay78}, TR-58.)

  \item[Relay~64G (stator earth fault):] High-penetration IBR generate
    sub-harmonic currents near $f_{\rm sub} \in [15, 25]\,\text{Hz}$ from
    converter dead-time and LCL-filter resonance.  When the conventional fixed
    injection at $f_{\rm conv} = 20\,\text{Hz}$ coincides with $f_{\rm sub}$,
    the noise-to-signal ratio at the relay input inverts: false alarms exceed
    99\% and missed detections occur for all high-resistance faults.
    (Section~\ref{sec:ch6_relay64g}, TR-61.)
\end{description}

All six relay-function extensions are complete as of April~2026.  Chapter~\ref{ch:system_integration}
integrates the full nine-element SAMBP suite at the system level.

% ──────────────────────────────────────────────────────────────────────────────
\section{Generator Differential for DFIG Machines (87G)}
\label{sec:ch6_87g}

\subsection{The DFIG Fault Current Challenge}
\label{sec:ch6_87g_problem}

DFIG wind turbines present a two-regime fault current that no existing relay
model captures.
While the rotor-side converter (RSC) remains active (\emph{pre-crowbar}
phase), the stator current behaves as a current-limited IBR:
\begin{equation}
  i_{\rm pre}(t) = k_{\rm ibr}\,\sin(\omega_0 t + \phi_{\rm pre}),
  \label{eq:ch6_dfig_pre}
\end{equation}
with no DC offset and no slip-frequency content.

After crowbar firing at $t = t_{\rm cb}$ the rotor circuit becomes a simple
$RL$ loop.  The analytical solution for the stator current introduces two
additional components:
\begin{equation}
  i_{\rm post}(t) =
    I_{\rm fund}\sin(\omega_0 t + \phi_{\rm fund})
    + I_{\rm nat}\,e^{-t_{\rm rel}/\tau_s}\sin(\omega_s\,t_{\rm rel} + \phi_{\rm nat})
    + I_{\rm DC}\,e^{-t_{\rm rel}/\tau_s},
  \label{eq:ch6_dfig_post}
\end{equation}
where $t_{\rm rel} = \max(t - t_{\rm cb}, 0)$, $\omega_s = s\,\omega_0$ is the
slip frequency ($\omega_s \in [9.4, 110]\,\text{rad/s}$ for $s \in [0.03, 0.35]$),
and $\tau_s = L_s'/(R_s + R_{\rm cb})$ is the stator transient time constant.

\begin{remark}[Physics-based collinearity reduction]
The DC aperiodic and natural-response components share the same time constant
$\tau_s$ because both originate from the same stator transient flux-linkage
initial condition.  This physical constraint reduces the parameter count from
11 to 10 and guarantees $\kappa_n \in [13, 67]$ for the 10-parameter model,
versus $\kappa_n \in [3, 36]$ for the SG model (Proposition~C9 applied).
\end{remark}

\subsection{Full 10-Parameter Piecewise-ODE Model}
\label{sec:ch6_dfig_model}

The complete DFIG model transitions smoothly between phases using a
$\sigma_h = 0.5\,\text{ms}$ smoothed Heaviside:
\begin{equation}
  i_{\rm diff}(t,\,\boldsymbol{\theta}) =
    H_r(t)\,i_{\rm pre}(t) + H_\sigma(t)\,i_{\rm post}(t,\,t_{\rm rel}),
  \label{eq:ch6_dfig_full}
\end{equation}
where $H_\sigma(t) = [1 + \exp(-(t - t_{\rm cb})/\sigma_h)]^{-1}$ and
$H_r = 1 - H_\sigma$.  The transition width (10\%--90\%) is $\approx 1.1\,\text{ms}$
— below the 1\,ms sample interval — making $\partial i_{\rm diff}/\partial t_{\rm cb}$
continuous everywhere (required for gradient-based optimisation).

The 10-parameter vector is
$\boldsymbol{\theta} = [k_{\rm ibr},\,\phi_{\rm pre},\,t_{\rm cb},\,I_{\rm fund},\,
\phi_{\rm fund},\,I_{\rm nat},\,\phi_{\rm nat},\,\omega_s,\,\tau_s,\,I_{\rm DC}]^\top$.

\subsection{CUSUM Crowbar Detector}
\label{sec:ch6_cusum}

The crowbar firing time $t_{\rm cb}$ must be detected online from the relay
current waveform because it is not observable from relay communications.
A CUSUM detector on the band-power of the slip-frequency component triggers when
the cumulative sum of slip-frequency amplitude increments exceeds a calibrated
threshold:
\begin{equation}
  g_k = \max(0,\; g_{k-1} + P_k - P_0 - K), \quad \text{alarm if } g_k \geq h,
  \label{eq:ch6_cusum}
\end{equation}
where $P_k$ is the short-time DFT band-power in $[\omega_s - \Delta\omega,\,
\omega_s + \Delta\omega]$ at sliding window $k$.
The Wald-optimal parameters are $K = \tfrac{1}{2}(P_{\rm post} - P_{\rm pre})
\approx 4.5\,P_0$ and $h = 1.5\,K$.
Two-cycle windows are used to reduce leakage from the 50\,Hz fundamental into
the slip band: main-lobe width $\approx 50\,\text{Hz}$ with Hanning weighting
gives $< 1\%$ leakage into $[1.5, 17.5]\,\text{Hz}$.

The CUSUM detector has three properties:
\begin{itemize}
  \item Detection latency $t_{95} \leq 40\,\text{ms}$ for slip $\geq 6\%$.
  \item False-positive rate $< 0.1\%$ for pure IBR current (no crowbar).
  \item Broadband slip coverage: all $|s| \in [0.03, 0.35]$ monitored with no
    a~priori slip estimate required.
\end{itemize}

\subsection{Two-Pass Levenberg--Marquardt Estimator}
\label{sec:ch6_dfig_lm}

A physics-initialised two-pass LM estimator fits $\boldsymbol{\theta}$:

\begin{description}
  \item[Pass 1 — Full window, all 10 parameters:]
    The initial guess places $t_{\rm cb}^{(0)}$ at the CUSUM alarm time,
    $k_{\rm ibr}^{(0)}$ from the pre-crowbar RMS, and
    $I_{\rm nat}^{(0)} = 0.3\,\text{pu}$, $\omega_s^{(0)} = 31.4\,\text{rad/s}$
    (5\% slip at 50\,Hz).
    Pass~1 solves all 10 parameters on the complete window
    (TRF trust-region solver, $\leq 2\,000$ function evaluations).

  \item[Pass 2 — Tail refinement of $\omega_s$ and $\tau_s$:]
    The final 3 cycles of the window (the ``tail'') contain primarily the
    decaying natural response.  Pass~2 fixes the 8 amplitude and phase
    parameters and refines only $(\omega_s, \tau_s)$ on the tail
    ($\leq 600$ evaluations), reducing $\kappa_n$ for the slow-parameter
    sub-problem from $\approx 67$ to $\approx 12$.
\end{description}

The condition number $\kappa_n \in [13, 67]$ across all slip groups tested,
remaining below $\kappa_{\max} = 80$.  At $s = 3\%$ ($\omega_s = 9.4\,\text{rad/s}$)
the slip frequency is closest to DC, creating the highest collinearity
($\kappa_n = 67$); Pass~2 tail refinement is critical in this regime.

\subsection{Validation}
\label{sec:ch6_87g_val}

Validation proceeded in three stages (TR-56):

\begin{description}
  \item[Deterministic 10-scenario sweep:]
    Internal 3PH, SLG, and HIF faults and external faults at three slip levels.
    10/10 PASS; relay correct in all cases including the minimum-slip ($s=3\%$)
    high-$\kappa_n$ case.

  \item[Monte Carlo ($7 \times 1\,000$ trials per slip group):]
    CCR\,=\,100\% across all seven slip groups
    ($s \in \{3\%,6\%,10\%,15\%,20\%,30\%,35\%\}$);
    $P_{FA} = 0$;
    CUSUM $t_{95} \leq 80\,\text{ms}$ for all groups.

  \item[ANDES WTDTA1 validation:]
    IEEE~14-bus extended with a DFIG wind farm (ANDES~v1.9 WTDTA1 device,
    Bus~8).  Three fault resistances $R_f \in \{0.01, 0.10, 0.50\}\,\text{pu}$,
    3 external + 3 internal cases.
    6/6 correct relay decisions; zero false trips on external faults;
    LM convergence confirmed; monotone $\hat{I}_{\rm sub}$ vs.\ $R_f$
    consistent with DFIG physics.
\end{description}

\begin{table}[h]
  \centering
  \caption{TR-56 ANDES WTDTA1 validation results — 6-case matrix.}
  \label{tab:ch6_dfig_andes}
  \begin{tabular}{llccccc}
    \toprule
    Case & Scenario & $R_f$ (pu) & Trip? & $t_{\rm trip}$ (ms) &
          $\hat{I}_{\rm sub}$ (pu) & CUSUM (ms) \\
    \midrule
    \texttt{ext\_R001} & External & 0.01 & \textbf{N} & --- & ---   & ---   \\
    \texttt{ext\_R010} & External & 0.10 & \textbf{N} & --- & ---   & ---   \\
    \texttt{ext\_R050} & External & 0.50 & \textbf{N} & --- & ---   & ---   \\
    \texttt{int\_R001} & Internal & 0.01 & \textbf{Y} & 950 & 0.735 & 55.0  \\
    \texttt{int\_R010} & Internal & 0.10 & \textbf{Y} & 950 & 0.511 & 58.1  \\
    \texttt{int\_R050} & Internal & 0.50 & \textbf{Y} & 950 & 0.395 & 152.6 \\
    \bottomrule
  \end{tabular}
\end{table}

\noindent
The monotone decrease of $\hat{I}_{\rm sub}$ with $R_f$ (0.735 $\to$ 0.395\,pu)
is physically correct: higher fault resistance reduces fault current magnitude.
The increasing CUSUM latency at $R_f = 0.50\,\text{pu}$ (152.6\,ms) arises
from the smaller slip-band power increment when $I_{\rm nat}$ is small.

\paragraph{ANDES simulator adequacy (TR-59).}
A 10-case ANDES positive-sequence sweep confirmed that ANDES is
\emph{not} adequate for validating the \texttt{sync\_oc} subtransient
identification algorithm: the confidence score $\hat{\gamma} \in [0.500, 0.582]$
across all cases, uniformly below the adaptation threshold $\gamma_{\rm th} = 0.70$.
The root cause is that positive-sequence ANDES enforces symmetric phasors,
suppressing the negative- and zero-sequence content that the subtransient
estimator uses for identifiability.
This is a correct gate behaviour: the confidence threshold correctly withholds
adaptation when input waveform quality is insufficient.
Full-EMT validation (PSCAD explicit flux equations) is required for
\texttt{sync\_oc}, and is deferred to the TR-67 HIL campaign
(Chapter~\ref{ch:system_integration}).

% ──────────────────────────────────────────────────────────────────────────────
\section{Frequency Protection for Zero-Inertia GFM Islands (Relay~81)}
\label{sec:ch6_relay81}

\subsection{Problem: Variable Virtual Inertia Time Constant}
\label{sec:ch6_relay81_problem}

In a GFM-dominated island, the inertial frequency response is governed by the
virtual-synchronous-machine droop filter:
\begin{equation}
  \hat{f}(s) = f_0 - \frac{\Delta P}{D_v} \cdot \frac{1}{1 + \tau_f s},
  \label{eq:ch6_gfm_freq}
\end{equation}
where $D_v$ is the droop coefficient, $\Delta P$ the power imbalance, and
$\tau_f \in [0.05, 1.0]\,\text{s}$ the virtual filter time constant — a
20$\times$ range across different GFM manufacturers.
Fixed ROCOF thresholds calibrated at one end of this range produce either
false trips (too-low $\tau_f$) or missed islanding events (too-high $\tau_f$).

\subsection{IFDI Estimator}
\label{sec:ch6_relay81_ifdi}

TR-63 proposes a 6-parameter GFM island model
$\{f_0, D_v, \tau_f, \Delta P, \Delta f_\infty, W_\tau\}$ estimated by a
two-pass Levenberg--Marquardt procedure with multi-start
$\tau_f \in \{0.05, 1.0\}\,\text{s}$.
The instantaneous frequency deviation index (IFDI) is:
\begin{equation}
  \mathrm{IFDI} = \frac{|\Delta\hat{f}_\infty|}{0.035} \times W_\tau,
  \label{eq:ch6_ifdi}
\end{equation}
where $\Delta\hat{f}_\infty$ is the estimated steady-state frequency deviation
and $W_\tau$ is a quality weight penalising uncertain $\tau_f$ estimates.
The trip threshold IFDI $\geq 1.15$ provides a 15\% security margin above the
unity normalisation point.

\subsection{Validation}
\label{sec:ch6_relay81_val}

Nine deterministic scenarios (three islanding events $\times$ three GFM
penetration levels) produced 9/9 PASS.
Monte Carlo ($N = 10\,000$ per hypothesis class) confirmed:
$P_D = 0.9947$, $P_{FA} = 0.0032$, all acceptance criteria met.
The non-detection zone (NDZ) for $|\Delta P| < 3\%$ produces $P_D = 0.377$
in the ambiguous region, which is expected and documented: the confidence gate
classifies events below 3\% as ambiguous and withholds adaptation, preserving
relay security while correctly identifying genuine islanding events at any
practically relevant power imbalance.

% ──────────────────────────────────────────────────────────────────────────────
\section{Solar PV Line Differential with Two-Parameter Zone Model (87L-PV)}
\label{sec:ch6_87l_pv}

\subsection{Problem: DC-Free PV Fault Current and Wasted Estimator Degrees of Freedom}
\label{sec:ch6_pv_problem}

A grid-connected PV inverter with a synchronous-reference-frame (SRF) current
controller drives $\mathbf{i}_{dq} \to \mathbf{i}_{dq}^*$ with bandwidth
$\omega_c \approx 300$--$800\,\text{rad/s}$.  This gives a DC-free per-phase
fault current:

\begin{proposition}[Zero DC Offset in PV Fault Current, TR-62 Proposition~1]
\label{prop:ch6_pv_zero_dc}
For a GFL inverter with SRF current controller of bandwidth $\omega_c$, the
per-phase fault current satisfies
\begin{equation}
  i(t) = k_{\rm ibr}\sin(\omega_0 t + \varphi) + O(e^{-\omega_c t}),
  \label{eq:ch6_pv_current}
\end{equation}
with the transient deviation decaying to below $1\%$ of $k_{\rm ibr}$ within
$t_s = 4.6/\omega_c \leq 15\,\text{ms}$ for $\omega_c \geq 300\,\text{rad/s}$.
\end{proposition}

The critical implication: the 4-parameter SAMBP SG model
$i_{\rm diff}(t) = I_{\rm fund}\sin(\omega t + \varphi) + I_{\rm DC}\,e^{-t/\tau_{\rm DC}}$
forces two of its four parameters ($I_{\rm DC}, \tau_{\rm DC}$) to fit noise,
causing collinearity ($\kappa_n \approx 36$ at a half-cycle window) and
delaying trip confirmation by 20--40\,ms.

\subsection{Two-Parameter PV Model}
\label{sec:ch6_pv_model}

The physically correct zone model for PV-fed lines is:
\begin{equation}
  \boxed{i_{\rm diff}^{(\rm PV)}(t) = I_{\rm fund}\sin(\omega t + \varphi), \qquad
  \boldsymbol{\theta}_{\rm PV} = [I_{\rm fund},\;\varphi] \in \mathbb{R}^2.}
  \label{eq:ch6_2param}
\end{equation}

\begin{proposition}[Unit Condition Number of 2-Parameter Model, TR-62 Proposition~2]
\label{prop:ch6_kappa1}
For a window spanning an integer number of periods $T_0$, the normalised
condition number of the analytical Jacobian of \eqref{eq:ch6_2param} is:
\begin{equation}
  \kappa_n(J_{\rm PV}) = 1.
  \label{eq:ch6_kappa1}
\end{equation}
\end{proposition}
\noindent Proof: $J_{\rm PV}^T J_{\rm PV} = (N/2)\,\text{diag}(1, I_f^2)$, which is diagonal;
column normalisation yields the identity matrix with $\sigma_{\max} = \sigma_{\min} = 1$.

The practical consequence is that the Stage-2 confidence gate accepts the
PV estimator from the first post-inception cycle.  The binding constraint
becomes $n_{\rm confirm} \geq 2$ cycles for security — giving a minimum trip
time of $20\,\text{ms}$ at 50\,Hz, versus $\geq 40\,\text{ms}$ for the
4-parameter SG path.

\begin{table}[h]
  \centering
  \caption{Condition number vs.\ window length for 4-parameter SG and
           2-parameter PV models.}
  \label{tab:ch6_kappa_window}
  \begin{tabular}{lcccc}
    \toprule
    Window & $N$ & $\kappa_n$ (4-param SG) & $\kappa_n$ (2-param PV) & Trip delay (SG / PV) \\
    \midrule
    0.5 cycle &  10 & 35.9 & 1.00 & 5 cyc / 1 cyc \\
    1.0 cycle &  20 &  5.4 & 1.00 & 3 cyc / 1 cyc \\
    2.0 cycles & 40 &  3.3 & 1.00 & 2 cyc / 1 cyc \\
    3.0 cycles & 60 &  2.9 & 1.00 & 1 cyc / 1 cyc \\
    \bottomrule
  \end{tabular}
\end{table}

\subsection{$k_{\rm ibr}$-Driven Model Selector}
\label{sec:ch6_pv_selector}

At each relay cycle, the model selector routes the estimation to either the
2-parameter PV path or the existing 4-parameter SG path using the
\textbf{DC decay ratio}:
\begin{equation}
  \mathrm{DDR} = \frac{I_{\rm DC}}{\max(I_{\rm fund},\,\varepsilon)},
  \qquad \varepsilon = 10^{-9},
  \label{eq:ch6_ddr}
\end{equation}
from a lightweight Pass-1 fit of the 4-parameter model.
Physical calibration: PV faults give DDR $\leq 0.05$; DFIG faults
DDR $\in [0.05, 0.20]$; SG faults DDR $\in [0.30, 1.50]$.
The PV model is selected when DDR $< 0.15$ (midway between PV and DFIG
upper limits, providing $3\times$ margin in each direction), subject to an
exponential veto: if $\Delta_{\rm peak}/I_{\rm fund} \geq 2.0$ (SG-signature
DC offset detectable in the first quarter-cycle), the selector returns to the
SG path regardless of DDR.

\subsection{Validation}
\label{sec:ch6_pv_val}

\begin{table}[h]
  \centering
  \caption{TR-62 Monte Carlo results — 2-parameter PV model ($N = 5\,000$ trials).}
  \label{tab:ch6_pv_mc}
  \begin{tabular}{lrrll}
    \toprule
    Metric & Value & Target & Status \\
    \midrule
    $P_D$ (detection, PV internal faults) & 0.9984 & $\geq 0.995$ & PASS \\
    $P_{FA}$ (false alarm, PV external) & 0.0000 & $\leq 0.001$ & PASS \\
    $t_{50}$ median trip time & 20\,ms & $\leq 20\,\text{ms}$ & PASS \\
    $t_{95}$ 95th-percentile trip time & 20\,ms & $\leq 40\,\text{ms}$ & PASS \\
    DDR accuracy (correct model routing) & 0.9984 & $\geq 0.98$  & PASS \\
    \bottomrule
  \end{tabular}
\end{table}

All five acceptance criteria are met.  The 8 non-detected trials (0.16\%)
represent edge cases where CT amplitude error combined with pre-fault load offset
pushes $\hat{I}_{\rm fund}$ below the Stage-2 threshold of $0.058\,\text{pu}$.
These are correctly classified as RESTRAIN rather than false alarms.

Three validation streams confirm the model consistently:
\begin{itemize}
  \item \textbf{Analytical} (TR-62 Propositions~1--2): $\kappa_n = 1$ proved;
    minimum trip time $20\,\text{ms}$ derived.
  \item \textbf{ANDES PVD1} (TR-62): IEEE~14-bus with 10$\times$PVD1 at Bus~4;
    6/6 correct relay decisions ($\kappa_n = 1.00$ measured in every internal
    case), zero false trips on external faults.
  \item \textbf{Python MC} (TR-62a): 9/9 deterministic PASS, all 5-metric MC
    targets met.
\end{itemize}

% ──────────────────────────────────────────────────────────────────────────────
\section{Loss of Excitation Protection with GFM Admittance Correction (Relay~40)}
\label{sec:ch6_relay40}

\subsection{Failure Modes under GFM-IBR Reactive Injection}
\label{sec:ch6_loe_fm}

Standard Relay~40 detects loss of excitation (LOE) by monitoring the imaginary
part of the generator admittance $Y_{\rm SG}(t) = G_{\rm sg} + jB_{\rm sg}(t)$
and tripping when $B_{\rm sg}$ crosses a threshold $B_{\rm trip}$ for a
persistence interval $T_{\rm delay}$.
In high-penetration GFM networks two failure modes arise:

\begin{description}
  \item[FM1 — Nuisance trip (security).]
    A GFM unit in reactive-\emph{absorption} mode contributes $jB_{\rm GFM} > 0$
    to the terminal bus.  The measured total admittance
    $Y_{\rm total} = Y_{\rm SG} + Y_{\rm GFM}$ may cross $B_{\rm trip}$ while
    the generator is healthy ($B_{\rm sg} < B_{\rm trip}$).

  \item[FM2 — Masked LOE (dependability).]
    A GFM unit in reactive-\emph{supply} mode contributes $jB_{\rm GFM} < 0$.
    The composite $\Im(Y_{\rm total}) = B_{\rm sg} - |B_{\rm GFM}|$ remains
    below $B_{\rm trip}$ even when $B_{\rm sg} > B_{\rm trip}$, concealing a
    genuine LOE event.
\end{description}

\subsection{GFM Admittance Correction}
\label{sec:ch6_loe_correction}

The correction applies an online GFM admittance estimate to the relay input:
\begin{equation}
  Y_{\rm corrected}(t) = Y_{\rm total}(t) - \widehat{Y}_{\rm GFM}(t),
  \label{eq:ch6_ycorr}
\end{equation}
where $\widehat{Y}_{\rm GFM}$ is obtained from PMU-grade measurements at the
GFM point of connection ($\sigma_\varepsilon \leq 0.03\,\text{pu}$).
The relay evaluates $B_{\rm trip}$ on $\Im(Y_{\rm corrected})$ rather than
$\Im(Y_{\rm total})$.

\begin{proposition}[TR-60 Proposition~1 — FM1 correction]
  \label{prop:ch6_fm1}
  Under reactive-absorption GFM operation, the corrected relay is secure
  (no false trip) provided the estimation error $|\varepsilon|$ satisfies
  $|\varepsilon| < B_{\rm trip} - B_{\rm sg,0}$.
\end{proposition}

\begin{proposition}[TR-60 Proposition~2 — FM2 correction]
  \label{prop:ch6_fm2}
  Under reactive-supply GFM operation masking a genuine LOE, the corrected
  relay trips provided $|\varepsilon| < B_{\rm sg} - B_{\rm trip}$.
\end{proposition}

Both conditions are satisfied with PMU accuracy
($3\sigma = 0.09\,\text{pu}$) given margins
$B_{\rm trip} - B_{\rm sg,0} = 0.25\,\text{pu}$ and
$B_{\rm sg} - B_{\rm trip} \approx 0.40\,\text{pu}$ (factor $> 3.4\times$
headroom for PMU errors).

\subsection{Validation Results}
\label{sec:ch6_loe_results}

TR-60 validates the correction on a 10-scenario matrix (S01--S05 security,
S06--S10 dependability) and a 2\,000-trial Monte Carlo.

\begin{table}[h]
  \centering
  \caption{TR-60 Monte Carlo results — Relay 40 LOE with GFM admittance correction.}
  \label{tab:ch6_loe_mc}
  \begin{tabular}{lcccc}
    \toprule
    Metric & Uncorrected & Corrected & Target & Status \\
    \midrule
    $P_D$ (dependability) & 0.568 & 1.000 & $\geq 0.970$ & PASS \\
    $P_{\rm FA}$ (false alarm) & 0.038 & 0.000 & $\leq 0.010$ & PASS \\
    Deterministic (10 scenarios) & 4/10 & 10/10 & 10/10 & PASS \\
    \bottomrule
  \end{tabular}
\end{table}

The $+43.2$\,pp improvement in dependability and elimination of false alarms
demonstrate that the GFM admittance correction fully restores Relay~40
performance in the presence of GFM reactive injection up to
$B_{\rm GFM,max} = 0.80\,\text{pu}$.
The $t_{50}$ trip time under LOE is $3\,786\,\text{ms}$ (governed by the
$T_{\rm delay} = 3\,\text{s}$ persistence interval, unchanged from the
conventional relay setting).

% ──────────────────────────────────────────────────────────────────────────────
\section{Out-of-Step Protection for IBR-Affected Generators (Relay~78)}
\label{sec:ch6_relay78}

\subsection{The IBR Synchronising Deficit Problem}
\label{sec:ch6_relay78_problem}

The classical equal-area criterion (EAC) determines the critical unstable
equilibrium point (CUEP) $\delta_u$ as the rotor angle at which the
decelerating area first equals the accelerating area during a fault.
For a single-machine infinite-bus (SMIB) system with no IBR,
$\delta_u = \pi - \delta_0$ where $\delta_0 = \arcsin(P_m/P_{\rm maxsg})$.
The conventional Relay~78 blinder is set at this angle.

When GFL IBR constitutes a fraction $k_{\rm gfl}$ of the total generation,
the post-fault synchronising power is reduced by the GFL active-power deficit
during low-voltage ride-through (LVRT):
\begin{equation}
  \Delta P_{\rm max,GFL} = \max\bigl(0,\;
    P_{\rm GFL,pre} - V_{\rm min} \cdot I_{\rm max,GFL} \cdot \cos\phi_{\rm LVRT}\bigr).
  \label{eq:ch6_gfl_deficit}
\end{equation}
The effective maximum post-fault power becomes
$P_{\rm maxeff} = P_{\rm maxsg} - \Delta P_{\rm max,GFL}$,
reducing $\delta_u$ below $\pi - \delta_0$.
The conventional blinder at $\pi - \delta_0$ therefore lies \emph{outside}
the true CUEP and is non-conservative.

\subsection{Conservative CUEP Lower Bound (Proposition~1)}
\label{sec:ch6_relay78_prop}

\begin{proposition}[Conservative CUEP Lower Bound, TR-58 Proposition~1]
\label{prop:ch6_cuep}
Let $P_{\rm maxeff} = P_{\rm maxsg} - \Delta P_{\rm max,GFL}$ as in
\eqref{eq:ch6_gfl_deficit}.  The closed-form lower bound
\begin{equation}
  \delta_{\rm LB} = \pi - \arcsin\!\left(\frac{P_m}{P_{\rm maxeff}}\right)
  \label{eq:ch6_delta_lb}
\end{equation}
satisfies $\delta_{\rm LB} \leq \delta_u$ for all physically realisable
GFL and GFM operating conditions.
\end{proposition}

\noindent Proof outline (TR-58 Appendix): the post-fault energy function
$V(\delta) = \int_{\delta_0}^{\delta}(P_m - P_e(\delta'))\,\mathrm{d}\delta'$
evaluated at $\delta_{\rm LB}$ is strictly positive (conservative) under the
GFL deficit model.
For GFM scenarios, the virtual synchronising torque $K_v(\delta - \delta_0)$
causes the true CUEP to approach $\pi$ (increasing conservatism gap);
$\delta_{\rm LB}$ remains conservative with $\varepsilon_{\rm max} = 30.25^\circ$.

The proposed blinder is set at $\delta_{\rm blinder} = \delta_{\rm LB} - 5^\circ$
(5° security margin).

\subsection{Validation}
\label{sec:ch6_relay78_val}

Ten deterministic scenarios (S01--S10) spanning pure-SG, GFM-dominant,
GFL-dominant, mixed, and SLG fault cases were simulated using the swing
equation with RK45 integration (TR-58).  All 10/10 PASS for the proposed
blinder.

The key geometric check (scenario S05, 30\% GFL, deep fault):
\begin{itemize}
  \item Standard blinder at $153.6^\circ$ lies \emph{outside} the true
    CUEP at $151.1^\circ$ — non-conservative ($+2.5^\circ$ error).
  \item Proposed blinder at $146.1^\circ$ lies \emph{inside} the CUEP
    — conservative ($\varepsilon = 5.0^\circ$).
\end{itemize}

Monte Carlo validation (2\,000 trials, clearly-stable / clearly-unstable
partition) confirmed:
\begin{itemize}
  \item $P_{\rm conserv} = P(\delta_{\rm LB} \leq \delta_u) = 1.000\,000$
    (zero violations).
  \item $\varepsilon_{\rm max} = 30.25^\circ$ (GFM scenarios, $K_v > 0$).
  \item $P_D = 1.000$, $P_{FA} = 0.000$, $t_{50} = 470\,\text{ms}$.
\end{itemize}

% ──────────────────────────────────────────────────────────────────────────────
\section{Stator Earth Fault with Optimal Sub-harmonic Injection (Relay~64G)}
\label{sec:ch6_relay64g}

\subsection{Problem: IBR Sub-harmonic Noise Corrupts Fixed 20\,Hz Injection}
\label{sec:ch6_64g_problem}

Relay~64G injects a sub-harmonic voltage at the machine neutral (conventional:
$f_{\rm conv} = 20\,\text{Hz}$) and measures the return current through the
stator-to-ground capacitance $C_G$ and fault resistance $R_f$.

Without fault, only a capacitive displacement current flows:
\begin{equation}
  I_{\rm healthy}(f) = \frac{V_{\rm inj} \cdot \omega C_G}{\sqrt{1 + (\omega C_G R_N)^2}},
  \label{eq:ch6_Ihealthy}
\end{equation}
which \emph{increases with} $f$ as $X_c = 1/(\omega C_G)$ decreases.

With a stator earth fault, the winding impedance from neutral to the fault point
is negligible at sub-harmonic frequencies, giving a \emph{frequency-independent}
fault current:
\begin{equation}
  I_{\rm fault}(R_f) = \frac{V_{\rm inj}}{R_N + R_f}.
  \label{eq:ch6_Ifault}
\end{equation}

High-penetration IBR generate sub-harmonic currents near
$f_{\rm sub} \in [15, 25]\,\text{Hz}$ from converter dead-time, LCL-filter
resonance, and GFM droop dynamics.  When $f_{\rm sub} \approx f_{\rm conv} = 20\,\text{Hz}$:
the noise standard deviation $\sigma_n(20\,\text{Hz})$ is maximal,
causing false alarms ($P_{FA} = 0.819$, measured for the conventional relay)
and missed high-resistance fault detections.

\subsection{Optimal Injection Frequency}
\label{sec:ch6_64g_fopt}

The optimal frequency $f^*$ maximises the balanced margin between security
and dependability:
\begin{equation}
  f^* = \arg\max_{f \in [20,40]\,\text{Hz}}
  \min\!\Bigl(
    \underbrace{I_{\rm trip} - I_{\rm healthy}(f) - 2\sigma_n(f)}_{\text{security margin}},\;
    \underbrace{I_{\rm fault}(R_f) - I_{\rm trip} - 2\sigma_n(f)}_{\text{dependability margin}}
  \Bigr).
  \label{eq:ch6_fopt}
\end{equation}

\begin{proposition}[TR-61 — Optimal frequency bounds]
  \label{prop:ch6_fopt}
  For $f_{\rm sub} \in [15, 25]\,\text{Hz}$ and $A_{\rm noise} > 0$, the optimal
  injection frequency satisfies $f^* \in (f_{\rm sub}, f_{\rm up})$ where
  $f_{\rm up}$ is defined by $I_{\rm healthy}(f_{\rm up}) = I_{\rm trip}$
  (security crossover, $\approx 35\,\text{Hz}$).  In practice $f^* \in [25, 36]\,\text{Hz}$.
\end{proposition}

\noindent
The key insight: $I_{\rm fault}$ is \emph{frequency-independent}
(Equation~\eqref{eq:ch6_Ifault}).  Selecting $f^* \gg f_{\rm sub}$ reduces
$\sigma_n(f^*)$ to near-zero without affecting the fault signal, yielding a
$50\times$ noise reduction for $f^* - f_{\rm sub} \geq 8\,\text{Hz}$.

\subsection{Validation Results}
\label{sec:ch6_64g_results}

TR-61 validates the adaptive injection scheme on a 10-scenario deterministic
matrix and a 2\,000-trial Monte Carlo.

\begin{table}[h]
  \centering
  \caption{TR-61 Monte Carlo results — Relay 64G optimal sub-harmonic injection.}
  \label{tab:ch6_64g_mc}
  \small
  \begin{tabular}{lcccc}
    \toprule
    Metric & Conv.\ (20\,Hz) & Optimal ($f^*$) & Target & Status \\
    \midrule
    $P_D$ (dependability) & 0.813 & 1.000 & $\geq 0.970$ & PASS \\
    $P_{\rm FA}$ (false alarm) & 0.819 & 0.000 & $\leq 0.010$ & PASS \\
    $f^*$ mean & 20.0\,Hz & 27.2\,Hz & $\in [25,36]\,\text{Hz}$ & PASS \\
    Deterministic (10 scenarios) & 0/10 & 10/10 & 10/10 & PASS \\
    \bottomrule
  \end{tabular}
\end{table}

The conventional relay at 20\,Hz fails catastrophically ($P_{\rm FA} = 0.819$)
when IBR noise peaks near the injection frequency.
The adaptive relay achieves perfect performance by injecting at
$f^* \approx 27\,\text{Hz}$ — consistently $\geq 8\,\text{Hz}$ above the
noise peak $f_{\rm sub}$, consistent with Proposition~\ref{prop:ch6_fopt}.
The noise reduction factor $\sigma_n(27\,\text{Hz})/\sigma_n(20\,\text{Hz})
\approx 1/50$ is the primary mechanism for both the security and dependability
improvements.

% ──────────────────────────────────────────────────────────────────────────────
\section{Chapter Summary}
\label{sec:ch6_summary}

This chapter extended the SAMBP framework to the six generator and PV line
protection functions.  All six are validated as of April~2026:

\begin{table}[h]
  \centering
  \caption{Generator-suite validation summary (all TRs complete April~2026).}
  \label{tab:ch6_summary}
  \small
  \begin{tabular}{llcccl}
    \toprule
    Section & Function & Det.\ PASS & $P_D$ & $P_{\rm FA}$ & Source \\
    \midrule
    \S\ref{sec:ch6_87g}     & 87G DFIG       & 10/10 & 1.000$^a$         & 0.000 & TR-56 \\
    \S\ref{sec:ch6_relay81} & Relay 81 IFDI  &  9/9  & 0.9947            & 0.0032 & TR-63 \\
    \S\ref{sec:ch6_87l_pv}  & 87L-PV 2-param &  9/9  & 0.9984            & 0.0000 & TR-62 \\
    \S\ref{sec:ch6_relay40} & Relay 40 LOE   & 10/10 & 1.000             & 0.000 & TR-60 \\
    \S\ref{sec:ch6_relay78} & Relay 78 EAC   & 10/10 & 1.000$^b$         & 0.000 & TR-58 \\
    \S\ref{sec:ch6_relay64g}& Relay 64G 64G  & 10/10 & 1.000             & 0.000 & TR-61 \\
    \bottomrule
    \multicolumn{6}{l}{\scriptsize $^a$CCR\,=\,100\% across all 7 slip groups (7\,000 MC trials).}\\
    \multicolumn{6}{l}{\scriptsize $^b$Analytical guarantee: Proposition~\ref{prop:ch6_cuep} proves $\delta_{\rm LB} \leq \delta_u$
    always; 0 violations / 2\,000 MC.}\\
  \end{tabular}
\end{table}

Three technical themes unify the six extensions:

\begin{enumerate}
  \item \textbf{Physics-based model reduction.}  Each zone model exploits a
    physical constraint to reduce parameters: the shared $\tau_s$ in the
    DFIG model (§\ref{sec:ch6_87g}), the absence of DC in PV current
    (§\ref{sec:ch6_87l_pv}), and the frequency-independence of $I_{\rm fault}$
    in Relay~64G (§\ref{sec:ch6_relay64g}).  In each case the reduction lowers
    $\kappa_n$ and either accelerates convergence or eliminates false-alarm risk.

  \item \textbf{Analytical conservatism guarantees.}  Relay~78 achieves a
    zero-violation conservatism guarantee (Proposition~\ref{prop:ch6_cuep})
    rather than a probabilistic performance statement.  This is the strongest
    form of SAMBP contribution: from Monte Carlo validation to formal proof.
    Relay~40 Propositions~\ref{prop:ch6_fm1}--\ref{prop:ch6_fm2} provide
    analogous explicit error-budget conditions on PMU accuracy.

  \item \textbf{Online compensation of IBR network state.}  Three relay
    functions require real-time IBR state awareness: Relay~40 uses PMU-measured
    $\widehat{Y}_{\rm GFM}$; Relay~78 uses $\Delta P_{\rm max,GFL}$ from the
    LVRT model; Relay~64G adapts $f^*$ online using the spectral estimate of
    $f_{\rm sub}$.  All three are compatible with the IEC~61850 GOOSE
    broadcast already deployed in Chapter~\ref{ch:system_integration}.
\end{enumerate}

Chapter~\ref{ch:system_integration} integrates all nine SAMBP protection
functions at the system level using the IEEE~39-bus test network and the
N-1 cross-trip coordination study (TR-66, complete April~2026).
